\documentclass[11pt,a4paper]{article}
\usepackage[utf8]{inputenc}
\usepackage[spanish]{babel}
\usepackage{amsmath, amssymb}
\usepackage{geometry}
\geometry{top=2.5cm, bottom=2.5cm, left=2.5cm, right=2.5cm}

\begin{document}

\section*{Piense, dialogue y comparta}

\begin{itemize}
    \item[\textbf{1.}] Usted trabaja haciendo entregas para una tienda de lácteos. En la parte de atrás de su camioneta hay una caja de huevos. En la tienda de lácteos se han agotado las cuerdas elásticas, por lo que la caja no está atada. Se le ha dicho que conduzca con cuidado porque el coeficiente de fricción estática entre el cajón y la cama del camión es $0.600$. No está preocupado, porque están viajando por una carretera que parece perfectamente recta. Debido a su confianza y falta de atención, su velocidad ha aumentado arriba de $45.0 \text{ mi/h}$. De repente, ve una curva adelante con un letrero de advertencia que dice: ``Peligro: curva no peraltada con radio de curvatura de $35.0 \text{ m}$.'' Usted está a $15.0 \text{ m}$ de donde empieza la curva. ¿Qué puede hacer para salvar los huevos: (i) tomar la curva a $45.0 \text{ mi/h}$, (ii) frenar hasta detenerse antes de entrar en la curva para pensar en esto, o (iii) reducir la velocidad para tomar la curva a una velocidad más lenta? Discuta estas opciones con su grupo y determine si hay un mejor camino a seguir.
    \vspace*{9cm}

    \item[\textbf{2.}] \textbf{ACTIVIDAD} Encuentre un video de YouTube que muestre el ciclo completo de un juego del parque de diversiones llamado el ``Rodeo''. En este juego, una persona se encuentra contra una pared en el borde de un disco que gira alrededor de un eje vertical. Cuando el disco alcanza su rapidez de funcionamiento, un brazo eleva el disco a un ángulo para que gire alrededor de un eje que es casi horizontal. Como resultado, la persona se mueve en la parte superior de un círculo vertical, aparentemente sin apoyo, pero no cae hacia abajo. Usando la altura de una persona típica en el juego, calcule el radio del disco, utilizando una imagen detenida del disco en su ángulo más alto. Inicie el video de nuevo y utilice el cronómetro de su teléfono celular para medir el periodo de rotación del disco. (a) A partir de esta información, calcule la aceleración centrípeta de una persona en la parte superior del recorrido. (b) ¿Cómo se compara esta aceleración con la debida a la gravedad? (c) ¿Por qué no se cae una persona en la parte superior hacia abajo?
    \vspace*{8cm}

    \item[\textbf{3.}] \textbf{ACTIVIDAD} Encuentre un video de YouTube que muestra el ciclo completo para un juego del parque de diversiones llamado ``Rotor''. En este juego, una persona se encuentra contra una pared en un cilindro que gira alrededor de un eje vertical. Cuando el cilindro alcanza su rapidez de funcionamiento, el piso se cae alejándose y las personas permanecen suspendidas en la pared. Usando la altura de una persona típica en el juego, calcule el radio del cilindro, utilizando una imagen detenida del juego. Empiece de nuevo el video y utilice el cronómetro de su teléfono celular para medir el periodo de rotación del cilindro. Con esta información, determine el coeficiente mínimo de fricción estática necesario entre la persona y la pared para mantener a la persona suspendida.
    \vspace*{8cm}
\end{itemize}

\newpage

\section*{Problemas}

\subsection*{SECCIÓN 6.1 Extensión del modelo de partícula en movimiento circular uniforme}

\begin{itemize}
    \item[\textbf{1.}] En el modelo de Bohr del átomo de hidrógeno, la rapidez del electrón es aproximadamente $2.20 \times 10^6 \text{ m/s}$. Encuentre (a) la fuerza que actúa sobre el electrón mientras da vueltas en una órbita circular de $0.529 \times 10^{-10} \text{ m}$ de radio y (b) la aceleración centrípeta del electrón.
    \vspace*{6cm}

    \item[\textbf{2.}] Mientras dos astronautas del \textit{Apolo} estaban en la superficie de la Luna, un tercer astronauta la orbitaba. Suponga que la órbita es circular y $100 \text{ km}$ arriba de la superficie de la Luna, donde la aceleración debida a la gravedad es $1.52 \text{ m/s}^2$. El radio de la Luna es $1.70 \times 10^6 \text{ m}$. Determine (a) la rapidez orbital del astronauta y (b) el periodo de la órbita.
    \vspace*{6cm}

    \item[\textbf{3.}] Un automóvil viaja inicialmente hacia el este y da vuelta al norte al viajar en una trayectoria circular con rapidez uniforme, como se muestra en la figura P6.3. La longitud del arco $ABC$ es $235 \text{ m}$ y el automóvil completa la vuelta en $36.0 \text{ s}$. (a) ¿Cuál es la aceleración cuando el automóvil está en $B$, ubicado a un ángulo de $35.0^\circ$? Exprese su respuesta en términos de los vectores unitarios $\hat{i}$ y $\hat{j}$. Determine (b) la rapidez promedio del automóvil y (c) su aceleración promedio durante el intervalo de $36.0 \text{ s}$.
    \vspace*{8cm}

    \item[\textbf{4.}] Una curva en una carretera forma parte de un círculo horizontal. Cuando la rapidez de un automóvil que circula por ella es de $14 \text{ m/s}$ constante, la fuerza total horizontal sobre el conductor tiene $130 \text{ N}$ de magnitud. ¿Cuál es la fuerza horizontal total sobre el conductor si la rapidez es $18.0 \text{ m/s}$?
    \vspace*{5cm}
\end{itemize}

\newpage

\begin{itemize}
    \item[\textbf{5.}] En un ciclotrón (un tipo de acelerador de partículas), un deuterón (de $2.00 \text{ u}$ de masa) alcanza una rapidez final de $10.0\%$ la rapidez de la luz mientras se mueve en una trayectoria circular de $0.480 \text{ m}$ de radio. ¿Qué magnitud de fuerza magnética se requiere para mantener el deuterón en una trayectoria circular?
    \vspace*{6cm}

    \item[\textbf{6.}] ¿Por qué no es posible la siguiente situación? El objeto de masa $m = 4.00 \text{ kg}$ en la figura P6.6 se une a una barra vertical mediante dos cuerdas de longitud $\ell = 2.00 \text{ m}$. Las cuerdas están unidas a la varilla en puntos separados por una distancia $d = 3.00 \text{ m}$. El objeto gira en un círculo horizontal con rapidez constante $v = 3.00 \text{ m/s}$, y las cuerdas permanecen tensas. La varilla gira junto con el objeto de forma que las cuerdas no se enrollan en la varilla. \textbf{¿Qué pasaría si?} ¿Esta situación sería posible en otro planeta?
    \vspace*{7cm}

    \item[\textbf{7.}] Usted está trabajando durante sus vacaciones de verano como un operador de juegos del parque de diversiones. El juego que está operando consiste en un cilindro vertical grande que gira sobre su eje lo suficientemente rápido para que cualquier persona en el interior de este se sostenga contra la pared cuando se retira el piso (figura P6.7). El coeficiente de fricción estática entre una persona de masa $m$ y la pared es $\mu_s$, y el radio del cilindro es $R$. Usted está girando el juego con una rapidez angular $\theta$ sugerida por su supervisor. (a) Suponga que una persona muy pesada sube al juego. ¿Necesita aumentar la rapidez angular para que esta persona no resbale por la pared? (b) Supongamos que alguien que sube al juego lleva un traje deportivo de satín muy resbaladizo. En este caso, ¿necesita aumentar la rapidez angular para que esta persona no resbale por la pared?
    \vspace*{7cm}

    \item[\textbf{8.}] Un conductor está demandando al Departamento Estatal de Carreteras después de un accidente en una autopista curva. El conductor perdió el control y se estrelló contra un árbol situado a poca distancia del borde exterior de la carretera. Él afirma que el radio de curvatura de la carretera no peraltada era demasiado pequeño para la rapidez límite, lo que causó que se deslizara hacia el exterior en la curva y pegara contra el árbol. Usted ha sido contratado como un testigo experto para la defensa, y se le ha pedido que use su conocimiento de la física para testificar que el radio de curvatura de la carretera es apropiado para el límite de rapidez. Los reglamentos estatales muestran que el radio de curvatura de una carretera peraltada en la que el límite de velocidad es de $65 \text{ mi/h}$ debe ser al menos $150 \text{ m}$. Usted construye un acelerómetro, que es una plomada con un transportador que usted cuelga en el techo de su auto. Un compañero que se sube a su vehículo observa que la plomada cuelga a un ángulo de $15.0^\circ$ de la vertical cuando el auto se conduce con la rapidez más segura de $23.0 \text{ m/s}$ en la curva en cuestión. ¿Cuál es su testimonio sobre el radio de la curva?
    \vspace*{7cm}
\end{itemize}

\newpage

\subsection*{SECCIÓN 6.2 Movimiento circular no uniforme}

\begin{itemize}
    \item[\textbf{9.}] Un halcón vuela en un arco horizontal de $12.0 \text{ m}$ de radio con una rapidez constante de $4.00 \text{ m/s}$. (a) Encuentre su aceleración centrípeta. (b) El halcón continúa volando a lo largo del mismo arco horizontal pero aumenta su rapidez a una proporción de $1.20 \text{ m/s}^2$. Encuentre la aceleración (magnitud y dirección) en esta situación en el momento que la rapidez del halcón es $4 \text{ m/s}$.
    \vspace*{6cm}

    \item[\textbf{10.}] Un niño de $40.0 \text{ kg}$ se mece en un columpio sostenido por dos cuerdas, cada una de $3.00 \text{ m}$ de largo. La tensión en cada cuerda en el punto más bajo es $350 \text{ N}$. Encuentre (a) la rapidez del niño en el punto más bajo y (b) la fuerza que ejerce el asiento sobre el niño en el punto más bajo. (Desprecie la masa del asiento.)
    \vspace*{6cm}

    \item[\textbf{11.}] Un niño de masa $m$ se mece en un columpio sostenido por dos cuerdas, cada una de longitud $R$. Si la tensión en cada cuerda en el punto más bajo es $T$, encuentre (a) la rapidez del niño en el punto más bajo y (b) la fuerza que ejerce el asiento sobre el niño en el punto más bajo. (Desprecie la masa del asiento.)
    \vspace*{6cm}

    \item[\textbf{12.}] Un extremo de una cuerda está fijo y un objeto pequeño de $0.500 \text{ kg}$ se ata al otro extremo, donde se balancea en una sección de un círculo vertical de $2.00 \text{ m}$ de radio, como se muestra en la figura P6.12. Cuando $\theta = 20.0^\circ$, la rapidez del objeto es $8.00 \text{ m/s}$. En este instante, encuentre (a) la tensión en la cuerda, (b) las componentes tangencial y radial de la aceleración y (c) la aceleración total. (d) ¿Su respuesta cambia si el objeto se balancea hacia su punto más bajo en lugar de balancearse hacia arriba? Explique su respuesta para el inciso (d).
    \vspace*{8cm}
\end{itemize}

\newpage

\begin{itemize}
    \item[\textbf{13.}] Una montaña rusa en el parque de diversiones \textit{Six Flags Great America} en Gurnee, Illinois, incorpora cierta tecnología de diseño ingeniosa y algo de física básica. Cada bucle vertical, en lugar de ser circular, tiene forma de lágrima (figura P6.13). Los carros viajan en el interior del bucle en la parte superior, y las magnitudes de velocidad son lo suficientemente grandes para asegurar que los carros permanezcan en la pista. El bucle más grande tiene $40.0 \text{ m}$ de alto. Suponga que la rapidez en la parte superior es $13.0 \text{ m/s}$ y la aceleración centrípeta correspondiente es $2g$. (a) ¿Cuál es el radio del arco de la lágrima en la parte superior? (b) Si la masa total de un carro más los pasajeros es $M$, ¿qué fuerza ejerce el riel sobre el carro en la parte superior? (c) Suponga que la montaña rusa tiene un bucle circular de $20.0 \text{ m}$ de radio. Si los carros tienen la misma rapidez, $13.0 \text{ m/s}$ en la parte superior, ¿cuál es la aceleración centrípeta en la parte superior? (d) Comente acerca de la fuerza normal en la parte superior en la situación descrita en el inciso (c) y de las ventajas de tener bucles con forma de lágrima.
    \vspace*{8cm}
\end{itemize}

\subsection*{SECCIÓN 6.3 Movimiento en marcos acelerados}

\begin{itemize}
    \item[\textbf{14.}] Un objeto de masa $m = 5.00 \text{ kg}$, unido a una balanza de resorte, descansa sobre una superficie horizontal sin fricción, como se muestra en la figura P6.14. La balanza de resorte, unida al extremo frontal de un vagón tiene una lectura de cero cuando el vagón está en reposo. (a) Determine la aceleración del vagón si la balanza de resorte tiene una lectura constante de $18.0 \text{ N}$ cuando el carro está en movimiento. (b) ¿Qué lectura constante mostrará la balanza si el vagón se mueve con velocidad constante? (c) Describa las fuerzas sobre el objeto como lo observa alguien en el vagón y alguien en reposo fuera del vagón.
    \vspace*{7cm}

    \item[\textbf{15.}] Una persona está de pie sobre una báscula en un elevador. Cuando el elevador parte, la báscula tiene una lectura constante de $591 \text{ N}$. Más tarde, cuando el elevador se detiene, la lectura de la báscula es $391 \text{ N}$. Suponga que la magnitud de la aceleración es la misma durante la partida y el frenado. Determine: (a) el peso de la persona, (b) la masa de la persona y (c) la aceleración del elevador.
    \vspace*{7cm}

    \item[\textbf{16.}] \textbf{Problema de repaso.} Una estudiante está de pie con su mochila en el piso cerca de ella, en un elevador que acelera continuamente hacia arriba con aceleración $a$. El ancho del elevador es $L$. La estudiante da a su mochila una patada rápida en $t = 0$ y le imparte una rapidez $v$ que la hace deslizar a través del piso del elevador. Al tiempo $t$, la mochila golpea la pared opuesta. Encuentre el coeficiente de fricción cinética $\mu_k$ entre la mochila y el piso del elevador.
    \vspace*{7cm}

    \item[\textbf{17.}] Un pequeño contenedor de agua se coloca sobre el carrusel dentro de un horno de microondas en un radio de $12.0 \text{ cm}$ desde el centro. El carrusel gira de manera uniforme y da una revolución cada $7.25 \text{ s}$. ¿Qué ángulo forma la superficie del agua con la horizontal?
    \vspace*{6cm}
\end{itemize}

\newpage

\subsection*{SECCIÓN 6.4 Movimiento en presencia de fuerzas resistivas}

\begin{itemize}
    \item[\textbf{18.}] La masa de un automóvil deportivo es $1\,200 \text{ kg}$. La forma del cuerpo es tal que el coeficiente de arrastre aerodinámico es $0.250$ y el área frontal es $2.20 \text{ m}^2$. Si ignora todas las otras fuentes de fricción, calcule la aceleración inicial que tiene el automóvil si ha viajado a $100 \text{ km/h}$ y cuando cambia a neutral y se le permite deslizarse.
    \vspace*{6cm}

    \item[\textbf{19.}] \textbf{Problema de repaso.} Una persona que lava ventanas jala un rodillo de goma por una ventana vertical muy alta. El rodillo tiene $160 \text{ g}$ de masa y está montado en el extremo de una barra ligera. El coeficiente de fricción cinética entre el rodillo y el vidrio seco es $0.900$. La persona lo presiona contra la ventana con una fuerza que tiene una componente horizontal de $4.00 \text{ N}$. (a) Si ella jala el rodillo por la ventana a velocidad constante, ¿qué componente de fuerza vertical debe ejercer? (b) La persona aumenta la componente de fuerza hacia abajo en $25.0\%$, pero todas las otras fuerzas permanecen iguales. Encuentre la aceleración del rodillo en esta situación. (c) Luego el rodillo se mueve en una porción húmeda de la ventana, donde su movimiento es resistido por una fuerza de arrastre de fluido $R$ proporcional a su velocidad de acuerdo con $R = -20.0 v$, donde $R$ está en newtons y $v$ está en metros por segundo. Encuentre la velocidad terminal a la que se aproxima el rodillo, si supone que la persona ejerce la misma fuerza descrita en el inciso (b).
    \vspace*{8cm}

    \item[\textbf{20.}] Un pequeño trozo de espuma de poliestireno, material de empaque, se suelta desde una altura de $2.00 \text{ m}$ sobre el suelo. Hasta que llega a su rapidez terminal, la magnitud de su aceleración se conoce mediante $a = g - Bv$. Después de caer $0.500 \text{ m}$, la espuma de estireno en efecto alcanza su rapidez terminal y después tarda $5.00 \text{ s}$ más en llegar al suelo. (a) ¿Cuál es el valor de la constante $B$? (b) ¿Cuál es la aceleración en $t = 0$? (c) ¿Cuál es la aceleración cuando la rapidez es $0.150 \text{ m/s}$?
    \vspace*{7cm}

    \item[\textbf{21.}] Una esfera pequeña de $3.00 \text{ g}$ de masa se libera desde el reposo en $t = 0$ desde un punto bajo la superficie de un fluido viscoso. Se observa que la rapidez terminal es $v_T = 2.00 \text{ cm/s}$. Encuentre (a) el valor de la constante $b$ en la ecuación 6.2, (b) el tiempo $t$ en el que la esfera alcanza $0.632 v_T$ y (c) el valor de la fuerza resistiva cuando la esfera alcanza su rapidez terminal.
    \vspace*{7cm}
\end{itemize}

\newpage

\begin{itemize}
    \item[\textbf{22.}] Suponga que la fuerza resistiva que actúa sobre un patinador rápido es $f = -kmv^2$, donde $k$ es una constante y $m$ es la masa del patinador. El patinador cruza la línea de meta de una competencia en línea recta con rapidez $v_i$ y después disminuye su rapidez deslizándose sobre sus patines. Demuestre que la rapidez del patinador en cualquier tiempo $t$ después de cruzar la línea final es $v(t) = v_i / (1 + ktv_i)$.
    \vspace*{6cm}

    \item[\textbf{23.}] Usted puede sentir una fuerza de arrastre de aire sobre su mano si estira el brazo por afuera de una ventana abierta en un automóvil que se mueve rápidamente. \textit{Nota:} No se ponga en peligro. ¿Cuál es el orden de magnitud de esta fuerza? En su solución, establezca las cantidades que mida o estime y sus valores.
    \vspace*{6cm}
\end{itemize}

\subsection*{PROBLEMAS ADICIONALES}

\begin{itemize}
    \item[\textbf{24.}] Un auto viaja en el sentido de las manecillas del reloj con rapidez constante alrededor de una sección circular de una carretera horizontal como se muestra en la vista aérea de la figura P6.24. Encuentre la dirección de su velocidad y aceleración en la posición \textcircled{A} y (b) en la posición \textcircled{B}.
    \vspace*{6cm}

    \item[\textbf{25.}] Se utiliza una cuerda bajo una tensión de $50.0 \text{ N}$ para dar de vueltas a una piedra en un círculo horizontal de radio $2.50 \text{ m} $ a una rapidez de $20.4 \text{ m/s}$ en una superficie sin fricción, como se muestra en la figura P6.25. Conforme se jala de la cuerda, aumenta la rapidez de la piedra. Cuando la cuerda sobre la mesa es de $1.00 \text{ m}$ de largo y la rapidez de la piedra es $51.0 \text{ m/s}$, la cuerda se rompe. ¿Cuál es la resistencia al rompimiento de la cuerda, en newtons?
    \vspace*{6cm}

    \item[\textbf{26.}] Perturbado por el exceso de velocidad de los automóviles fuera de su lugar de trabajo, el premio Nobel, Arthur Holly Compton, diseñó un tope (llamado el ``tope Holly'') e lo instaló. Suponga que un auto de $1\,800 \text{ kg}$ pasa por un tope en un camino que sigue un arco de un círculo de radio $20.4 \text{ m}$ como se muestra en la figura P6.26. (a) Si el auto viaja a $30.0 \text{ km/h}$, ¿qué fuerza ejerce la carretera sobre el auto cuando el automóvil pasa por el punto más alto del tope? (b) \textbf{¿Qué pasaría si?} ¿Cuál es la rapidez máxima que el automóvil puede tener sin perder contacto con la carretera cuando pasa este punto más alto?
    \vspace*{7cm}

    \item[\textbf{27.}] Un automóvil de masa $m$ pasa sobre un tope en un camino que sigue el arco de un círculo de radio $R$, como se muestra en la figura P6.26. (a) ¿Qué fuerza ejerce el camino sobre el automóvil mientras este pasa el punto más alto del tope si viaja a una rapidez $v$? (b) \textbf{¿Qué pasaría si?} ¿Cuál es la máxima rapidez que puede tener el automóvil mientras pasa el punto más alto sin perder contacto con el camino?
    \vspace*{6cm}
\end{itemize}

\newpage

\begin{itemize}
    \item[\textbf{28.}] El juguete de un niño consiste en una pequeña cuña que tiene un ángulo agudo $\theta$ (figura P6.28). El lado inclinado de la cuña no tiene fricción, y un objeto de masa $m$ en él se mantiene a una altura constante si la cuña se hace girar a cierta rapidez constante. La cuña está girando al rotar su eje, una varilla vertical que está firmemente sujeta a la cuña en el extremo inferior. Demuestre que, cuando el objeto se encuentra en reposo en un punto a una distancia $L$ arriba a lo largo de la cuña, la rapidez del objeto debe ser $v = (gL \operatorname{sen}\theta)^{1/2}$.
    \vspace*{6cm}

    \item[\textbf{29.}] Un hidroavión de masa total $m$ acuatiza en un lago con rapidez inicial $v_i\hat{i}$. La única fuerza horizontal sobre este es una fuerza resistiva sobre sus flotadores desde el agua. La fuerza resistiva es proporcional a la velocidad del hidroavión: $\vec{R} = -b\vec{v}$. La segunda ley de Newton aplicada al hidroavión es $-bv\hat{i} = m(dv/dt)\hat{i}$. A partir del teorema fundamental del cálculo, esta ecuación diferencial implica que la relación de cambio es de acuerdo con
    $$ \int_{v_i}^{v} \frac{dv}{v} = -\frac{b}{m} \int_{0}^{t} dt $$
    (a) Realice las integraciones para determinar la rapidez del hidroavión como función del tiempo. (b) Trace una gráfica de la rapidez como función del tiempo. (c) ¿El hidroavión se detiene por completo después de un intervalo de tiempo finito? (d) ¿El hidroavión viaja una distancia finita para detenerse?
    \vspace*{8cm}

    \item[\textbf{30.}] Un objeto de masa $m_1 = 4.00 \text{ kg}$ está atado a un objeto de masa $m_2 = 3.00 \text{ kg}$ con la cuerda 1 de longitud $\ell = 0.500 \text{ m}$. La combinación se hace oscilar en una trayectoria circular vertical sobre una segunda cuerda, la cuerda 2, de longitud $\ell = 0.500 \text{ m}$. Durante el movimiento, las dos cuerdas son colineales en todo momento como se muestra en la figura P6.30. En la parte superior de su movimiento, $m_2$ viaja a $v = 4.00 \text{ m/s}$. (a) ¿Cuál es la tensión en la cuerda 1 en este instante? (b) ¿Cuál es la tensión en la cuerda 2 en este instante? (c) ¿Cuál cuerda se romperá primero si la combinación se gira más y más rápido?
    \vspace*{7cm}

    \item[\textbf{31.}] Una bola de masa $m = 0.275 \text{ kg}$ se mantiene en una trayectoria circular vertical en una cuerda de longitud $L = 0.850 \text{ m}$ como en la figura P6.31. (a) ¿Cuáles son las fuerzas que actúan en la bola en cualquier punto de la trayectoria? (b) Dibuje diagramas de fuerza para la pelota cuando se encuentra en la parte inferior del círculo y cuando está en la cima. (c) Si su rapidez es de $5.20 \text{ m/s}$ en la parte superior del círculo, ¿cuál es la tensión en la cuerda ahí? (d) si la cuerda se rompe cuando su tensión excede $22.5 \text{ N}$, ¿cuál es la rapidez máxima que puede tener la pelota en el fondo antes de que ocurra?
    \vspace*{8cm}
\end{itemize}

\newpage

\begin{itemize}
    \item[\textbf{32.}] ¿Por qué no es posible la siguiente situación? Un niño travieso va a un parque de diversiones con su familia. En un juego, después de una severa reprimenda de su madre, él se desliza fuera de su asiento y trepa a la parte superior de la estructura del juego, que tiene forma de cono con su eje vertical y sus lados inclinados, haciendo un ángulo de $\theta = 20.0^\circ$ con la horizontal como se muestra en la figura P6.32. Esta parte de la estructura gira alrededor del eje central vertical cuando opera el juego. El niño se sienta en la superficie inclinada en un punto $d = 5.32 \text{ m}$ debajo del centro del cono por el lado oblicuo del cono y está enfadado. El coeficiente de fricción estática entre el niño y el cono es $0.700$. El operador del juego no se da cuenta de que el niño está lejos de su asiento y así continúa operando el juego. Como resultado, el niño sentado, haciendo pucheros gira en una trayectoria circular a una rapidez de $3.75 \text{ m/s}$.
    \vspace*{7cm}

    \item[\textbf{33.}] El piloto de un avión ejecuta una maniobra de rizo con rapidez constante en un círculo vertical. La rapidez del avión es $300 \text{ mi/h}$ en la parte superior del rizo y $450 \text{ mi/h}$ en la parte inferior y el radio del círculo es $1\,200 \text{ pies}$. (a) ¿Cuál es el peso aparente del piloto en el punto más bajo si su peso verdadero es $160 \text{ lb}$? (b) ¿Cuál es su peso aparente en el punto más alto? (c) \textbf{¿Qué pasaría si?} Describa cómo experimentaría el piloto la sensación de ausencia de peso si puede variar el radio y la rapidez. \textit{Nota:} Su peso aparente es igual a la magnitud de la fuerza que ejerce el asiento sobre su cuerpo.
    \vspace*{7cm}

    \item[\textbf{34.}] Una vasija que rodea un drenaje tiene la forma de un cono circular que se abre hacia arriba, y en todas partes tiene un ángulo de $35.0^\circ$ con la horizontal. Un cubo de hielo de $25.0 \text{ g}$ se hace deslizar alrededor del cono sin fricción en un círculo horizontal de radio $R$. (a) Encuentre la rapidez que debe tener el cubo de hielo como función de $R$. (b) ¿Es necesaria alguna parte de la información para la solución? Suponga que $R$ se hace dos veces más grande. (c) ¿La rapidez requerida aumenta, disminuye o permanece constante? Si cambia, ¿en qué factor? (d) ¿El tiempo requerido para cada revolución aumenta, disminuye o permanece constante? Si cambia, ¿en qué factor? (e) ¿Las respuestas a los incisos (c) y (d) parecen contradictorias? Explique.
    \vspace*{8cm}

    \item[\textbf{35.}] \textbf{Problema de repaso.} Mientras aprende a conducir, usted está en un automóvil de $1\,200 \text{ kg}$ que se mueve a $20.0 \text{ m/s}$ a través de un gran estacionamiento vacío y a nivel. Súbitamente se da cuenta de que se dirige justo hacia una pared de ladrillos de un gran supermercado y está en peligro de chocar con ella. El pavimento puede ejercer una fuerza horizontal máxima de $7\,000 \text{ N}$ en el automóvil. (a) Explique por qué debería esperar que la fuerza tenga un valor máximo bien definido. (b) Suponga que pisa los frenos y no gira el volante. Encuentre la distancia mínima a la que debe estar de la pared para evitar un choque. (c) Si no frena y en vez de ello mantiene rapidez constante y gira el volante, ¿cuál es la distancia mínima a la que debe estar de la pared para evitar un choque? (d) ¿Cuál método, (b) o (c), es mejor para evitar una colisión? O, ¿debe usar tanto frenos como volante, o ninguno? Explique. (e) ¿La conclusión del inciso (d) depende de los valores numéricos que se proporcionan en este problema, o es verdad en general? Explique.
    \vspace*{8cm}
\end{itemize}

\newpage

\begin{itemize}
    \item[\textbf{36.}] Un camión se mueve con aceleración constante $a$ en una colina que hace un ángulo $\phi$ con la horizontal como en la figura P6.36. Una pequeña esfera de masa $m$ está suspendida desde el techo de la camioneta por un cable ligero. Si el péndulo hace un ángulo constante $\theta$ con la perpendicular al techo, ¿a qué es igual $a$?
    \vspace*{6cm}

    \item[\textbf{37.}] Ya que la Tierra gira en torno a su eje, un punto sobre el ecuador experimenta una aceleración centrípeta de $0.033\ 7 \text{ m/s}^2$, mientras que un punto en los polos no experimenta aceleración centrípeta. Si una persona en el ecuador tiene una masa de $75.0 \text{ kg}$, calcule (a) la fuerza gravitacional (peso real) sobre la persona y (b) la fuerza normal (peso aparente) de la persona. (c) ¿Cuál es mayor? Suponga que la Tierra es una esfera uniforme y considere $g = 9.800 \text{ m/s}^2$.
    \vspace*{7cm}

    \item[\textbf{38.}] Un disco de masa $m_1$ se une a una cuerda y se le permite dar vueltas en un círculo de radio $R$ sobre una mesa horizontal sin fricción. El otro extremo de la cuerda pasa a través de un pequeño orificio en el centro de la mesa, y una carga de masa $m_2$ se une a la cuerda (figura P6.38). La carga suspendida permanece en equilibrio mientras que el disco en la mesa da vueltas. Encuentre expresiones simbólicas para (a) la tensión en la cuerda, (b) la fuerza radial que actúa sobre el disco y (c) la rapidez del disco? (d) Describa cualitativamente qué ocurrirá en el movimiento del disco si el valor de $m_2$ aumenta un poco al colocar una carga adicional sobre este. (e) Describa cualitativamente qué ocurrirá en el movimiento del disco si el valor de $m_2$ disminuye al remover una parte de la carga suspendida.
    \vspace*{8cm}

    \item[\textbf{39.}] Galileo pensó acerca de si la aceleración debía definirse como la relación de cambio de la velocidad en el tiempo o como la relación de cambio en velocidad en la distancia. Él eligió la anterior, así que use el nombre ``vroomosidad'' para la función de cambio de la velocidad con la distancia. Para el movimiento de una partícula en una línea recta con aceleración constante, la ecuación $v = v_i + at$ da su velocidad $v$ como función del tiempo. De igual modo, para el movimiento lineal de una partícula con vroomosidad constante $k$, la ecuación $v = v_i + kx$ da la velocidad como función de la posición $x$ si la rapidez de la partícula es $v_i$ en $x = 0$. (a) Encuentre la ley que describe la fuerza total que actúa en este objeto, de masa $m$. (b) Describa un ejemplo de tal movimiento o explique por qué tal movimiento es irreal. Considere (c) la posibilidad de $k$ positiva y también (d) la posibilidad de $k$ negativa.
    \vspace*{8cm}
\end{itemize}

\newpage

\begin{itemize}
    \item[\textbf{40.}] A los integrantes de un club de paracaidismo se les dieron los siguientes datos para usar en la planeación de sus saltos. En la tabla, $d$ es la distancia que cae desde el reposo un paracaidista en una ``posición extendida estable en caída libre'' en función del tiempo de caída $t$. (a) Convierta las distancias en pies a metros. (b) Grafique $d$ (en metros) en función de $t$. (c) Determine el valor de la rapidez terminal $v_T$ al encontrar la pendiente de la porción recta de la curva. Aplique un ajuste de mínimos cuadrados para determinar esta pendiente.

    \begin{center}
    \begin{tabular}{cc|cc|cc}
    \hline
    $t \text{ (s)}$ & $d \text{ (ft)}$ & $t \text{ (s)}$ & $d \text{ (ft)}$ & $t \text{ (s)}$ & $d \text{ (ft)}$ \\
    \hline
    0 & 0 & 7 & 652 & 14 & 1 831 \\
    1 & 16 & 8 & 808 & 15 & 2 005 \\
    2 & 62 & 9 & 971 & 16 & 2 179 \\
    3 & 138 & 10 & 1 138 & 17 & 2 353 \\
    4 & 242 & 11 & 1 309 & 18 & 2 527 \\
    5 & 366 & 12 & 1 483 & 19 & 2 701 \\
    6 & 504 & 13 & 1 657 & 20 & 2 875 \\
    \hline
    \end{tabular}
    \end{center}
    \vspace*{8cm}

    \item[\textbf{41.}] Un automóvil recorre una curva peraltada como se explicó en el ejemplo 6.4 y se muestra en la figura 6.5. El radio de curvatura del camino es $R$, el ángulo de peralte es $\theta$ y el coeficiente de fricción estática es $\mu_s$. (a) Determine el intervalo de rapidez que puede tener el automóvil sin deslizarse arriba o abajo del peralte. (b) Encuentre el valor mínimo para $\mu_s$ tal que la rapidez mínima sea cero.
    \vspace*{7cm}

    \item[\textbf{42.}] En el ejemplo 6.5, investigamos las fuerzas que un niño experimenta en una rueda de la fortuna. Suponga que los datos en ese ejemplo se aplican a este problema. ¿Qué fuerza (magnitud y dirección) ejerce el asiento sobre un niño de $40.0 \text{ kg}$ cuando este está a medio camino entre la parte superior e inferior?
    \vspace*{6cm}

    \item[\textbf{43.}] \textbf{Problema de repaso.} Una porción de masilla inicialmente se ubica en el punto $A$ en el borde de una rueda de molino que gira en torno a un eje horizontal. La masilla se desplaza del punto $A$ cuando el diámetro a través de $A$ es horizontal. Luego se eleva verticalmente y regresa a $A$ en el instante en que la rueda completa una revolución. Usando esta información, se quiere encontrar la rapidez $v$ de la masilla, cuando sale de la rueda y la fuerza que la sostenía a la rueda. (a) ¿Qué modelo de análisis es apropiado para el movimiento de la masilla conforme se levanta y cae? (b) Use este modelo para encontrar una expresión simbólica para el intervalo de tiempo entre que la masilla deja el punto $A$ y cuando llega a $A$, en términos de $v$ y $g$. (c) ¿Cuál es el modelo de análisis adecuado para describir un punto sobre la rueda? (d) Encuentre el periodo del movimiento del punto $A$ en términos de la rapidez tangencial $v$ y el radio $R$ de la rueda. (e) Haga el intervalo de tiempo del inciso (b) igual al periodo del inciso (d) y resuelva para la rapidez $v$ de la masilla cuando deja la rueda. (f) Si la masa de la masilla es $m$, ¿cuál es la magnitud de la fuerza que la mantuvo en la rueda antes de ser liberada?
    \vspace*{8cm}

    \item[\textbf{44.}] Un modelo de avión de masa $0.750 \text{ kg}$ vuela con una rapidez de $35.0 \text{ m/s}$ en un círculo horizontal en el extremo de un cable de control de $60.0 \text{ m}$ de largo como se muestra en la figura P6.44a. En la figura P6.44b se muestran las fuerzas ejercidas sobre el avión: la tensión en el alambre de control, la fuerza gravitacional y la elevación aerodinámica que actúa en $\theta = 20.0^\circ$ hacia adentro de la vertical. Calcule la tensión en el cable, suponiendo que hace un ángulo constante de $\theta = 20.0^\circ$ con la horizontal.
    \vspace*{7cm}
\end{itemize}

\newpage

\subsection*{PROBLEMAS DE DESAFÍO}

\begin{itemize}
    \item[\textbf{45.}] Un objeto de $9.00 \text{ kg}$ que parte del reposo cae a través de un medio viscoso y experimenta una fuerza resistiva dada por la ecuación 6.2. El objeto alcanza la mitad de su rapidez terminal en $5.54 \text{ s}$. (a) Determine la rapidez terminal. (b) ¿En qué tiempo la rapidez del objeto es tres cuartos de la rapidez terminal? (c) ¿Qué distancia recorrió el objeto en los primeros $5.54 \text{ s}$ de movimiento?
    \vspace*{8cm}

    \item[\textbf{46.}] Para $t < 0$, un objeto de masa $m$ no experimenta fuerza y se mueve en la dirección $x$ positiva con una rapidez constante $v_i$. Comenzando en $t = 0$, cuando el objeto pasa de la posición $x = 0$, experimenta una fuerza neta resistiva proporcional al cuadrado de su rapidez: $\vec{F}_{\text{neta}} = -mkv^2 \hat{i}$, donde $k$ es una constante. La rapidez del objeto después de $t = 0$ viene dada por $v = v_i / (1 + kv_it)$. (a) Encuentre la posición $x$ del objeto como una función del tiempo. (b) Encuentre la velocidad del objeto como una función de la posición.
    \vspace*{7cm}

    \item[\textbf{47.}] Un golfista saca desde una posición precisamente a $\phi_i = 35.0^\circ$ latitud norte. Golpea la bola hacia el sur, con un alcance de $285 \text{ m}$. La velocidad inicial de la bola está a $48.0^\circ$ sobre la horizontal. Suponga que la resistencia del aire es despreciable para la pelota de golf. (a) ¿Cuánto tiempo está la bola en vuelo? El hoyo está hacia el sur de la posición del golfista, y haría un hoyo en uno si la Tierra no girara. La rotación de la Tierra hace que el punto de partida se mueva en un círculo de radio $R_E \cos\phi_i = (6.37 \times 10^6 \text{ m})\cos 35.0^\circ$, como se muestra en la figura P6.47. El punto de partida completa una revolución cada día. (b) Encuentre la rapidez hacia el este del punto de partida, en relación con las estrellas. El hoyo también se mueve al este, pero está $285 \text{ m}$ más al sur y por tanto a una latitud ligeramente menor $\phi_f$. Dado que el hoyo se mueve en un círculo ligeramente más grande, su rapidez debe ser mayor que la del punto de partida. (c) ¿Por cuánto la rapidez del hoyo supera la del punto de partida? Durante el intervalo de tiempo en que la bola está en vuelo, se mueve arriba y abajo así como al sur con el movimiento de proyectil que estudió en el capítulo 4, pero también se mueve al este con la rapidez que encontró en el inciso (b). Sin embargo, el hoyo se mueve al este a una rapidez mayor, y jala la bola con la rapidez relativa que encontró en el inciso (c). (d) ¿A qué distancia hacia el oeste del hoyo aterriza la bola?
    \vspace*{9cm}
\end{itemize}

\newpage

\begin{itemize}
    \item[\textbf{48.}] Una sola cuenta puede deslizarse con fricción despreciable en un alambre rígido que se dobló en una espira circular de $15.0 \text{ cm}$ de radio, como se muestra en la figura P6.48. El círculo siempre está en un plano vertical y gira de manera estable en torno a su diámetro vertical con un periodo de $0.450 \text{ s}$. La posición de la cuenta se describe mediante el ángulo $\theta$ que la línea radial, desde el centro de la espira a la cuenta, forma con la vertical. (a) ¿A qué ángulo arriba del fondo del círculo puede permanecer la cuenta sin movimiento en relación con el círculo que gira? (b) \textbf{¿Qué pasaría si?} Repita el problema y considere que el periodo de rotación del círculo es $0.850 \text{ s}$. (c) Describa cómo la solución al inciso (b) es fundamentalmente diferente de la solución al inciso (a). (d) Para cualquier periodo o tamaño de espira, ¿siempre hay un ángulo al que la cuenta puede permanecer quieta en relación con la espira? (e) ¿Alguna vez hay más de dos ángulos? Arnold Arons sugirió la idea para este problema.
    \vspace*{9cm}

    \item[\textbf{49.}] Debido a la rotación de la Tierra, una plomada no cuelga exactamente a lo largo de una línea recta dirigida hacia el centro de la Tierra. ¿Cuánto se desvía la plomada de una línea radial a $35.0^\circ$ latitud norte? Suponga que la Tierra es esférica.
    \vspace*{6cm}

    \item[\textbf{50.}] Usted tiene un gran empleo: ¡trabajar en un estadio de béisbol de las grandes ligas durante el verano! En este estadio, la rapidez de cada lanzamiento se mide usando una pistola de radar dirigida al lanzador por un operador detrás de la base. El operador tiene tanta experiencia con este trabajo que ha perfeccionado una técnica por la cual puede hacer cada medida en el instante exacto en que la pelota sale de la mano del lanzador. Su supervisor le pide que construya un algoritmo que proporcione la rapidez con la que la bola cruza la base, a $18.3 \text{ m}$ del lanzador, con base en la rapidez medida $v_i$ de la pelota cuando sale de la mano del lanzador. La rapidez en la base será más baja debido a la fuerza resistiva del aire sobre la pelota de béisbol. El movimiento vertical de la bola es pequeño, por lo que, para una buena aproximación, podemos considerar solo el movimiento horizontal de la pelota. Usted empieza a desarrollar su algoritmo aplicando el modelo de partícula bajo una fuerza neta a la pelota de béisbol en la dirección horizontal. Mide que un lanzamiento tiene una rapidez de $40.2 \text{ m/s}$ cuando sale de la mano del jugador. Usted debe decirle a su supervisor la rapidez con la que viajaba, cuando cruzó la base. (\textit{Sugerencia:} Utilice la regla de la cadena para expresar la aceleración en términos de una derivada con respecto a $x$, y luego resuelva una ecuación diferencial para $v$ para encontrar una expresión para la rapidez de la pelota de béisbol en función de su posición. La función implicará una exponencial. También utilice la tabla 6.1.)
    \vspace*{8cm}
\end{itemize}

\end{document}
